\chapter{Kiểm thử}
\label{ch:test}

\section{Chiến lược kiểm thử}
Hệ thống được kiểm thử tự động ở ba mức, chạy bằng một lệnh \texttt{npm test}:
\begin{itemize}
  \item \textbf{Kiểm thử đơn vị} (\texttt{tests/rules.test.js}, 39 ca): kiểm tra các hàm thuần của bộ quy tắc nghiệp vụ, tập trung vào \emph{giá trị biên}: 9/10 học viên (online), 19/20 và 30/31 (offline), 39/40 (chia đều), các mốc quy đổi điểm CEFR.
  \item \textbf{Kiểm thử tích hợp API} (\texttt{tests/api.test.js}, 31 ca): khởi động máy chủ thật trên cổng ngẫu nhiên với \emph{CSDL tạm riêng} được nạp dữ liệu mẫu, gọi HTTP như trình duyệt, đi hết luồng đăng ký → thanh toán → thi → xếp lớp → tờ trình → phê duyệt → lập lịch → mở/huỷ lớp, đồng thời kiểm tra phân quyền.
  \item \textbf{Kiểm thử thời gian thực} (\texttt{tests/chat.test.js}, 8 ca): các client Socket.IO (một khách, hai tư vấn viên, một giáo vụ) kiểm tra hàng đợi, nhận phiên tranh chấp, nhắn tin hai chiều, khôi phục phiên, chặn socket lạ, đóng phiên.
\end{itemize}
Mỗi tệp kiểm thử dùng CSDL riêng trong thư mục tạm nên kết quả lặp lại được và không ảnh hưởng dữ liệu demo.

\section{Các ca kiểm thử tiêu biểu}
\begin{xltabular}{\linewidth}{|l|L{5.4cm}|Y|c|}
  \caption{Ca kiểm thử tiêu biểu và kết quả}\label{tab:tc}\\ \hline
  \headrow \hd{Mã} & \hd{Kịch bản} & \hd{Kết quả mong đợi} & \hd{KQ} \\ \hline
  \endfirsthead \hline
  \headrow \hd{Mã} & \hd{Kịch bản} & \hd{Kết quả mong đợi} & \hd{KQ} \\ \hline
  \endhead
  TC-01 & Lớp Online 9 / 10 học viên & BELOW\_MIN (thiếu 1) / ELIGIBLE & Đạt \\ \hline
  TC-02 & Lớp Offline 19, 20, 30, 31 học viên & BELOW\_MIN, ELIGIBLE, FULL, OVER\_CAPACITY & Đạt \\ \hline
  TC-03 & Xếp học viên thứ 30 / thứ 31 vào lớp Offline & ENROLLED / NEEDS\_APPROVAL & Đạt \\ \hline
  TC-04 & Xếp vào lớp Online đã đủ 40 & REJECT – ONLINE\_FULL, không lập tờ trình & Đạt \\ \hline
  TC-05 & Tờ trình 30 + 3, phòng 40 chỗ & Đề xuất KEEP; SPLIT dạng REMAINDER [30, 3] & Đạt \\ \hline
  TC-06 & Tờ trình 30 + 3, phòng 30 chỗ & KEEP không khả thi, đề xuất SPLIT & Đạt \\ \hline
  TC-07 & Tờ trình 30 + 10 & SPLIT dạng EVEN [20, 20] & Đạt \\ \hline
  TC-08 & Trùng phòng / trùng giảng viên cùng thứ – ca – thời gian & Phát hiện xung đột ROOM / TEACHER & Đạt \\ \hline
  TC-09 & Đăng ký hợp lệ & 201, mã \texttt{DK-YYYY-NNNN}, PENDING\_PAYMENT & Đạt \\ \hline
  TC-10 & Đăng ký trùng khoá; dữ liệu sai & 409; 400 kèm lỗi từng trường & Đạt \\ \hline
  TC-11 & Tra cứu hồ sơ đúng mã, sai SĐT & 404 & Đạt \\ \hline
  TC-12 & Làm bài khi chưa nộp học phí & 403 PAYMENT\_REQUIRED (BR-01) & Đạt \\ \hline
  TC-13 & Thanh toán 2 lần & Lần 2 trả 409 & Đạt \\ \hline
  TC-14 & Làm bài: trả lời đúng 15/20 & 20 câu không lộ đáp án; điểm 15, trình độ B2; hồ sơ TESTED & Đạt \\ \hline
  TC-15 & Gọi API quản trị không token; tư vấn viên xếp lớp & 401; 403 & Đạt \\ \hline
  TC-16 & Tham số tra cứu \verb|' OR 1=1 --| & 200, danh sách rỗng, không lỗi SQL & Đạt \\ \hline
  TC-17 & Giáo vụ ra quyết định tờ trình & 403 – chỉ Lãnh đạo & Đạt \\ \hline
  TC-18 & Lãnh đạo giữ ngoại lệ (hội trường 45 chỗ) & Lớp có 31 học viên, không còn hồ sơ chờ & Đạt \\ \hline
  TC-19 & Lãnh đạo tách lớp 30 + 12 = 42 & Hai lớp 21 + 21; lớp mới có \texttt{parent\_class\_id}, được gợi ý phòng trống & Đạt \\ \hline
  TC-20 & Từ chối tờ trình & Hồ sơ chờ trở về TESTED & Đạt \\ \hline
  TC-21 & Mở lớp Offline 21 HV chưa có phòng & 422 – “chưa được phân phòng học” & Đạt \\ \hline
  TC-22 & Gán phòng đang bận cùng ca & 422 SCHEDULE\_CONFLICT & Đạt \\ \hline
  TC-23 & Gán phòng trống rồi mở lớp & IN\_PROGRESS & Đạt \\ \hline
  TC-24 & Huỷ lớp & Toàn bộ học viên trở về danh sách chờ xếp lớp & Đạt \\ \hline
  TC-25 & Hai tư vấn viên cùng nhận một phiên & Người thứ hai bị từ chối & Đạt \\ \hline
  TC-26 & Tin nhắn hai chiều realtime & Bên kia nhận đúng nội dung, đúng người gửi & Đạt \\ \hline
  TC-27 & Socket lạ gửi tin vào phiên người khác & Bị từ chối & Đạt \\ \hline
  TC-28 & Đóng phiên rồi gửi tiếp & Bị từ chối; tin nhắn đã lưu CSDL & Đạt \\ \hline
  TC-29 & Tách lớp khi 30 + 35 hồ sơ chờ & Ba lớp 22 / 22 / 21 (không lớp nào vượt 30); hai lớp mới khác giảng viên, không trùng phòng & Đạt \\ \hline
  TC-30 & Sai mật khẩu 8 lần liên tiếp & Lần thứ 9 trả 429 kể cả khi nhập đúng & Đạt \\ \hline
  TC-31 & Giáo vụ nhận phiên chat & Bị từ chối (chỉ tư vấn viên/lãnh đạo) & Đạt \\ \hline
\end{xltabular}

\section{Kết quả chạy kiểm thử}
\begin{figure}[H]
\begin{PlainCode}
> node --test --test-reporter=spec "tests/*.test.js"
> BR-04 – phương án xử lý tờ trình
  v 30 + 3 = 33 < 40, phòng 40 chỗ -> đề xuất giữ ngoại lệ, tách phần dư
  v 30 + 3 nhưng phòng chỉ 30 chỗ -> không giữ được, đề xuất tách
  v 30 + 10 = 40 -> chia đều 20/20 (cả 2 lớp đủ điều kiện mở)
  ...
> Xếp lớp & tờ trình vượt sĩ số (BR-04)
  v Học viên thứ 30 vào lớp offline -> ENROLLED; thứ 31 -> chờ duyệt + tờ trình
  v Lãnh đạo "tách lớp" 30 + 12 = 42 -> chia đều 21 / 21, lớp mới có lịch + phòng
  ...
i tests 78
i suites 13
i pass 78
i fail 0
\end{PlainCode}
\captionof{lstlisting}{Kết quả \texttt{npm test} (rút gọn)}\label{lst:testout}
\end{figure}

\begin{table}[H]\centering
\caption{Tổng hợp kết quả kiểm thử tự động}\label{tab:testsum}
\begin{tabular}{|l|c|c|c|}
  \hline
  \headrow \hd{Tệp kiểm thử} & \hd{Số ca} & \hd{Đạt} & \hd{Thời gian} \\ \hline
  \texttt{rules.test.js} – quy tắc nghiệp vụ & 39 & 39 & $\approx$ 0,1 s \\ \hline
  \texttt{api.test.js} – tích hợp REST API & 31 & 31 & $\approx$ 1,5 s \\ \hline
  \texttt{chat.test.js} – Socket.IO & 8 & 8 & $\approx$ 1,1 s \\ \hline
  \textbf{Tổng} & \textbf{78} & \textbf{78 (100\%)} & \\ \hline
\end{tabular}
\end{table}

\section{Ma trận truy vết yêu cầu}
Bảng~\ref{tab:trace} nối từng yêu cầu của đề bài tới use case, mô-đun cài đặt và ca kiểm thử tương ứng, bảo đảm không yêu cầu nào bị bỏ sót và mọi quy tắc nghiệp vụ đều có kiểm thử.

\begin{xltabular}{\linewidth}{|L{2.7cm}|L{1.5cm}|Y|L{3.4cm}|}
  \caption{Truy vết yêu cầu – use case – mã nguồn – kiểm thử}\label{tab:trace}\\ \hline
  \headrow \hd{Yêu cầu} & \hd{Use case} & \hd{Mô-đun cài đặt} & \hd{Kiểm thử} \\ \hline
  \endfirsthead \hline
  \headrow \hd{Yêu cầu} & \hd{Use case} & \hd{Mô-đun cài đặt} & \hd{Kiểm thử} \\ \hline
  \endhead
  FR-01, FR-02 (CN1) & UC-01 & \texttt{services/catalog.js}, \texttt{public/js/portal.js} & API danh mục khoá học \\ \hline
  FR-03 (CN3) & UC-02 & \texttt{catalog.searchClasses} & API lọc lớp, TC-16 \\ \hline
  FR-04 (CN2) & UC-03 & \texttt{registration.create}, \texttt{rules.validateRegistration} & TC-09, TC-10 \\ \hline
  FR-05, FR-06 & UC-04 & \texttt{registration.pay}, \texttt{findByCodeAndPhone} & TC-11, TC-13 \\ \hline
  FR-07, BR-01 & UC-05 & \texttt{services/placement.js}, \texttt{rules.canTakePlacement} & TC-12, TC-14 \\ \hline
  FR-08 – FR-10 (CN4) & UC-06 & \texttt{sockets/chat.js}, \texttt{services/chat.js} & TC-25 – TC-28, TC-31 \\ \hline
  FR-11, BR-06 & UC-07 & \texttt{enrollment.enroll}, \texttt{rules.placementDecision} & TC-03, TC-04 \\ \hline
  FR-13, BR-04 & UC-08 & \texttt{enrollment.decide}, \texttt{rules.proposeResolution} & TC-05 – TC-07, TC-17 – TC-20, TC-29 \\ \hline
  FR-12, BR-05, BR-08 & UC-09 & \texttt{enrollment.updateArrangement}, \texttt{openClass}, \texttt{rules.findConflicts}, \texttt{canOpenClass} & TC-08, TC-21 – TC-24 \\ \hline
  BR-02, BR-03, BR-07 & – & \texttt{rules.evaluateClass}, view \texttt{v\_class\_enrollment} & TC-01, TC-02 \\ \hline
  FR-14 & UC-10 & \texttt{services/stats.js}, \texttt{public/js/admin.js} & API dashboard \\ \hline
  Bảo mật, phân quyền & – & \texttt{middleware/auth.js}, \texttt{middleware/throttle.js} & TC-15, TC-16, TC-30 \\ \hline
\end{xltabular}

\section{Kiểm thử giao diện}
Ngoài kiểm thử tự động, toàn bộ giao diện được chạy thử trên trình duyệt Chrome (độ phân giải 1440×900 và giả lập điện thoại 390×844) bằng công cụ \texttt{tools/screenshots.js}: công cụ tự đăng nhập từng vai trò, đi qua các màn hình, gõ tin nhắn chat ở cả hai phía và chụp ảnh – các ảnh trong Chương~\ref{ch:impl} được tạo ra từ kịch bản này. Quá trình chạy thử đã giúp phát hiện và sửa các lỗi: đường dẫn thư viện Chart.js bị chặn bởi trường \texttt{exports} của gói; xung đột tên lớp CSS giữa khối người dùng và bong bóng tin nhắn; lớp đang học bị gợi ý ưu tiên cho học viên mới (đã bổ sung trọng số giảm cho lớp đã khai giảng).
